% =========================================================================
% §7 Discussion (+ Conclusion)
% Budget: ~1.0 pp
% Drafting order: After Results + Integrity.
%
% Content blueprint:
%   - What the ablation says: index moves localization a lot
%     (+39.6 pp paired Wilcoxon, p<0.0001), moves resolve with
%     statistical separation (+7.9 pp, p=0.003), cost-neutral per cell
%     and favorable on $/solved. Causal because everything else held
%     identical across three seeds.
%   - Cross-harness validity: SC-ON matches or modestly favors OpenCode
%     on resolve (p=0.087) and View B acc@5 (p=0.080) at no cost penalty.
%     Reframe: structural codebase index is not duplicating agentic grep
%     behavior; at minimum it does not regress the agent.
%   - Where the index does NOT help: cross-harness comparisons are
%     marginal at conventional thresholds, directionally in SC-ON's favor
%     across seeds. The within-harness ablation is the cleaner read.
%   - Implications for harness design: when does a structural codebase
%     index pay off (multi-file changes), when doesn't it.
%   - Conclusion paragraph closes the doc.
%
% Constraints:
%   - Conclusion is ONE short paragraph; not a separate section.
%   - Stay scientific. Marketing reuses these numbers separately.
% =========================================================================

\section{Discussion}
\label{sec:discussion}

\textbf{What the ablation says.} SC-ON and SC-OFF differ only in the two engine tools
(\S\ref{sec:design-arms}, \S\ref{sec:system-engine}); model, prompts, sandbox, scorer,
seeds, caps, and the remaining tool set are held identical, so the within-harness deltas
read causally on the index. View~B acc@5 moves from 44.3\% to 84.5\% (paired Wilcoxon
$p<0.0001$), and resolve moves from 41.9\% to 50.4\% (paired $p=0.003$); per-cell mean
cost is statistically null ($p=0.73$) while turns and tokens both fall significantly. The
substantive read is that the index moves localization a lot, moves resolve significantly,
and is cost-neutral at the cell level while favorable on \$/solved.

\textbf{Cross-harness validity.} SC-ON matches or modestly favors OpenCode on resolve
($p=0.087$) and on View~B acc@5 ($p=0.080$), directionally in SC-ON's favor across seeds
but not at conventional significance. The structural codebase index does not duplicate
behavior that competent agentic grep already reaches; at minimum, against a competent
grep-and-read comparator, the index does not regress the agent on the metrics that decide
the task. The cross-harness effort gap is where the two arms separate cleanly: SC-ON
converges in fewer turns and fewer tokens (both $p<0.0001$, magnitudes in
Table~\ref{tab:headline}).

\textbf{Where the index does not help.} The ON vs.\ OpenCode comparisons on resolve and
View~B acc@5 are marginal at conventional significance thresholds (\S\ref{sec:results-crossharness});
the within-harness ablation is the cleanest read of what the index actually contributes.

\textbf{Implications for harness design.} We treat the implications as conditional
observations, not prescriptions. The largest causal localization gain in our data sits in
the 3-or-more-file gold bucket (\S\ref{sec:results-heterogeneity},
Figure~\ref{fig:filecount}), consistent with the call-graph intuition: when the change
spans files, a structural index that ranks paths by reachability pays off more than
agentic grep over the working copy. The cost-favorable finding simplifies the deployment
question: at a fixed model and comparable caps, a structural codebase index is the lower
\$/solved arm on these benchmarks, so the per-task win comes from localization quality
and the downstream turn savings it enables.

\paragraph{Conclusion.} This study reports a leak-audited, model-controlled, causal
ablation of a shipped structural codebase index inside a coding-agent harness, paired
with a cross-harness validity check against an agentic-grep comparator. The index
causally moves localization substantially and resolve with statistical separation
(50.4\% vs.\ 41.9\% on resolve, paired Wilcoxon $p=0.003$), at no per-cell cost penalty
and lower \$/solved than either the within-harness OFF arm or the OpenCode comparator. The released artifacts (the
exclusion ledger, the audit script, and the dual-view localization extractor) back the
integrity claims in \S\ref{sec:integrity} and the result numbers in \S\ref{sec:results}.
At this model and on these benchmarks, the deployment question for a structural codebase
index is not whether it is too expensive to run, but whether the workload includes
multi-file changes where structural ranking pays off.
